% Part III-B, abstract (included by main_en.tex).
When a verifier can label only some of a GRPO run's samples, a judge can fill in rewards for the rest. We study this
on MATH with Qwen3.5-4B-Base, using Part II's setup: the verifier labels 2 of the 8 completions of each group, and the
other 6 are dropped (\texttt{rlvr\_sub}), rewarded with \kev{}'s frozen judgment (\texttt{sub\_kev}), rewarded with
\kev{} after its pointer head is refitted online on the labels collected so far (\texttt{sub\_kev\_recal}), or rewarded
with a head of the same shape trained from scratch on the same features and labels (\texttt{sub\_probe}); two seeds
each, 200 steps. The frozen judge ranks the unlabelled completions well (AUROC \NUM{fillauc_all_subkev}) but is too
generous, and both of its policies learned to repeat themselves after the answer, which the judge pays for; against dropping
the completions they end \NUM{d_subkev_math} \ci{\NUM{dlo_subkev_math}}{\NUM{dhi_subkev_math}} points on MATH-500.
Refitting the head fixes its calibration (expected calibration error \NUM{fillece_all_subrecal} against
\NUM{fillece_all_subkev}), but one of its two policies still repeats itself. The head trained from scratch ends
\NUM{d_subprobe_math} \ci{\NUM{dlo_subprobe_math}}{\NUM{dhi_subprobe_math}} points against dropping on MATH-500 and
\NUM{d_subprobe_aime24} on AIME24, above the full-label baseline too. Both of its policies switched to opening their
solutions with an informal draft in plain text, a switch that in this project's other GRPO runs marks the best runs
whatever their reward, so two seeds cannot tell how much of the gain is the probe's. Code, logs, every training completion and every number in
this paper are in the repository.
